% Generated draft: requires mathematical and source review.
\documentclass[11pt]{article}
\usepackage[margin=25mm]{geometry}
\usepackage{fontspec}
\setmainfont{Latin Modern Roman}
\usepackage{amsmath,amssymb,mathtools,unicode-math}
\setmathfont{Latin Modern Math}
\usepackage{enumitem,needspace,xurl,hyperref}
\setlength{\parindent}{0pt}
\setlength{\parskip}{5pt}
\setlength{\emergencystretch}{3em}
\newcommand{\N}{\mathbb{N}}
\newcommand{\Z}{\mathbb{Z}}
\newcommand{\Q}{\mathbb{Q}}
\newcommand{\R}{\mathbb{R}}
\newcommand{\C}{\mathbb{C}}
\newcommand{\F}{\mathbb{F}}
\newcommand{\A}{\mathbb{A}}
\newcommand{\PP}{\mathbb{P}}
\newcommand{\E}{\mathbb{E}}
\newcommand{\eps}{\varepsilon}
\newcommand{\dd}{\,\mathrm{d}}
\begin{document}
\section*{UnsolvedMath ID: 1200014}
\textbf{Source problem code:} catalogue entry
\subsection{Definitions and mathematical statement}
The source asks for a quantitative version of an asserted compactness consequence: for every \(\varepsilon>0\), every finite graph can be approximated to error \(\varepsilon\) by a graph of bounded size \(K(\varepsilon)\), independently of the graph being approximated. The record does not define approximation error or graph size, and it does not specify a graph metric, normalization, or degree restrictions.

Conditionally, if \(\mathcal G\) is the intended class of finite graphs, \(d\) its intended approximation pseudometric, and size means vertex count, the bound sought would be an explicit finite function \(K:(0,\infty)\to\N\) such that
\[
\forall\varepsilon>0\ \forall G\in\mathcal G\ \exists H\in\mathcal G:
\quad |V(H)|\le K(\varepsilon),\qquad d(G,H)<\varepsilon.
\]
One could ask to estimate the optimal envelope
\[
K_{\mathcal G,d}(\varepsilon)=
\sup_{G\in\mathcal G}\min\{|V(H)|:H\in\mathcal G,\ d(G,H)<\varepsilon\},
\]
with an empty feasible set assigned value \(\infty\). This is only a conditional formulation: neither \(\mathcal G\), \(d\), nor the interpretation of size is supplied. The literal phrase ``of size \(K\)'' also does not distinguish exact size from an upper bound. Different graph-limit topologies or error measures produce different quantitative questions, so none is silently selected.

The supplied research summary labels the question OPEN-TRIAGE, describes existence as essentially compactness, and says no explicit estimate is known to its author and an effective version appears unresolved. That wording is a limited supplied assessment, not a verified claim that no estimate exists.
\subsection{Short English statement}
Estimate how large a finite graph must be to approximate any finite graph to a prescribed error. The approximation notion, graph class, and meaning of size are missing, so a definite quantitative problem cannot yet be extracted.
\subsection{Sources}
\begin{itemize}
\item Local UnsolvedMath registry, permanent ID \(1200014\), exact source code \texttt{catalogue entry}, ``Some Questions --- Question 14.''
\item Abert, \emph{Some questions} (2010), Question 14. Supplied author-hosted source: \url{https://www.renyi.hu/~abert/questions.pdf}; manual transcription; recorded access 2026-07-29.
\item Supplied OPEN-TRIAGE research summary, without a separate authoritative citation for the compactness theorem, its metric setting, or current quantitative bounds.
\end{itemize}
\end{document}
